\documentclass[leqno, a4paper, 12pt]{jsarticle}
\usepackage{amssymb}
\usepackage{amsthm}
\usepackage{amsmath}
\usepackage{comment}
\usepackage{url}

\usepackage{enumitem}
%\usepackage{showkeys}


%定理環境 thmはあまり使わずpropをなるべく使う。
%主張は基本的に証明の中で使い、そのため番号を付けない。
%注意環境を付けてみた。
\newtheorem{thm}{定理}
\newtheorem{prop}[thm]{命題}
\newtheorem{cor}[thm]{系}
\newtheorem{lem}[thm]{補題}
\newtheorem{df}[thm]{定義}
\newtheorem{exam}[thm]{例}
\newtheorem{fact}[thm]{事実}
\newtheorem*{claim}{主張}
\newtheorem*{cau}{注意}
\newtheorem*{rmk}{注意}
\renewcommand{\proofname}{証明}
%矢印たち。
%\toは\rightarrowと同じ。\getsは\leftarrowと同じ。同値は\iff
\newcommand{\To}{\Longrightarrow}
\newcommand{\Gets}{\LongLeftarrow}
%黒板太字
\newcommand{\nn}{\mathbb{N}}
\newcommand{\zz}{\mathbb{Z}}
\newcommand{\qq}{\mathbb{Q}}
\newcommand{\rr}{\mathbb{R}}
\newcommand{\cc}{\mathbb{C}}
%無限大
\newcommand{\mgn}{\infty}
%差集合は\setminusで統一する。等号の否定は\neq
%真部分集合は\subsetneqq　偏微分は\partial
%ディスプレイスタイル。\dfracでディスプレイ分数
\newcommand{\dpst}{\displaystyle}

\title{論文メモ
\large{\\--- ティーチェ・ウリゾーン？ ---}
}
\author{ナンブ　キトラ}
\date{\today}
\begin{document}
\maketitle

本棚を整理したいので，
溜まっている論文のメモを書いて未来への手紙とすることにする．


論文
\cite{MR0293605}
では空間
$X$がfunctionally Hausdorff 
であるための必要十分条件が
$X$の
任意の可算部分集合$X_{0}$
について連続関数
$f\colon X\to [0, 1]$
であって
$f|X_{0}$
が単射になるものが存在することだということを
証明している．
他にも正規性と$\aleph_{0}$-collectionwise normality
が同値であることを主張しているように見える．

論文
\cite{MR4431927}
では距離空間の上のティーチェの拡張定理の具体的な
やり方が紹介されてて面白い．
読むといいと思う．
本題としては拡張定理の同時化である．

%READ!
%myplain is the same to amsplain style. 
\nocite{*}
\bibliographystyle{myplaindoidoi}
\bibliography{TU.bib}



\end{document}